% Mathathon paper.
% Build: latexmk -pdf paper.tex   (or upload to Overleaf; it compiles with pdfLaTeX)
%
% The proof below follows Sections 1--8 of the team's notes
% "Why zeta(3) is irrational: a proof built from its requirements"
% (Peer Review/Zeta3-4-Motivated-proof-Codex/zeta3-motivated-proof.pdf),
% which develop Beukers' 1979 proof. Those notes have been checked by AI only.
\documentclass[11pt]{article}

\usepackage[a4paper,margin=25mm]{geometry}
\usepackage[T1]{fontenc}
\usepackage{lmodern}
\usepackage{amsmath,amssymb,amsthm}
\usepackage{graphicx}
\usepackage{booktabs}
\usepackage[colorlinks=true,linkcolor=blue,citecolor=blue,urlcolor=blue]{hyperref}
\usepackage[capitalise,nameinlink]{cleveref}

% ---------------------------------------------------------------- macros
\newcommand{\Z}{\mathbb{Z}}
\newcommand{\Q}{\mathbb{Q}}
\newcommand{\N}{\mathbb{N}}
\newcommand{\dd}{\,\mathrm{d}}
\newcommand{\Dn}{D_n}
\DeclareMathOperator{\lcm}{lcm}

% ---------------------------------------------------------------- theorems
\theoremstyle{plain}
\newtheorem{theorem}{Theorem}[section]
\newtheorem{lemma}[theorem]{Lemma}
\newtheorem{proposition}[theorem]{Proposition}
\newtheorem{corollary}[theorem]{Corollary}
\theoremstyle{definition}
\newtheorem{definition}[theorem]{Definition}
\newtheorem{example}[theorem]{Example}
\theoremstyle{remark}
\newtheorem{remark}[theorem]{Remark}

\title{Why \texorpdfstring{$\zeta(3)$}{zeta(3)} is irrational: Beukers' proof}
\author{Mathathon team}
\date{\today}

\begin{document}
\maketitle

\begin{abstract}
We give a complete, self-contained proof that $\zeta(3)=\sum_{m\ge1}m^{-3}$ is irrational
(Ap\'ery's theorem), following the integral method of Beukers. A double integral against the
kernel $-\log(xy)/(1-xy)$ turns every integer polynomial of degree $n$ into a linear form
$A\zeta(3)-B$ whose denominators are cleared by $\Dn^3$, where $\Dn=\lcm(1,\dots,n)$. Choosing
the shifted Legendre polynomial makes the form positive and smaller than
$2\zeta(3)(\sqrt2-1)^{4n}$. A short identity between $N!$ and least common multiples shows that
$\Dn<3^n$ for infinitely many $n$, which is enough to conclude; the prime number theorem is not
needed.
\end{abstract}

\tableofcontents

% ======================================================================
\section{The strategy: an integer between 0 and 1}\label{sec:strategy}

Suppose $\theta=p/q$ with $p\in\Z$ and $q\in\N$. For integers $A,B$ we have
$q(A\theta-B)=Ap-qB\in\Z$, so if $A\theta-B\ne0$ then $|A\theta-B|\ge1/q$. Since $q$ is fixed, a
sequence of \emph{nonzero} integer linear forms $A_n\theta-B_n$ tending to $0$ proves that $\theta$
is irrational.

We shall construct real numbers $J_n$ with three properties:
\begin{align}
  &\Dn^3J_n\in\Z+\Z\zeta(3), \qquad \Dn=\lcm(1,2,\dots,n)\quad(D_0=1), \label{eq:arith}\\
  &J_n>0, \label{eq:pos}\\
  &J_n\le C\rho^n \quad\text{with } \rho \text{ small enough.} \label{eq:small}
\end{align}

\begin{remark}
The obvious candidates, the partial sums $s_n=\sum_{m\le n}m^{-3}$, have the right denominators
but are far too poor approximations: $\zeta(3)-s_n\ge\int_{n+1}^\infty x^{-3}\dd x=\frac1{2(n+1)^2}$,
so $\Dn^3(\zeta(3)-s_n)\ge n^3/(2(n+1)^2)\to\infty$. Much more cancellation is needed, with the
arithmetic kept under control.
\end{remark}

% ======================================================================
\section{A kernel with the right arithmetic}\label{sec:kernel}

\subsection{Writing \texorpdfstring{$\zeta(3)$}{zeta(3)} as an integral}

Integration turns powers into reciprocals of integers:
\[
  \int_0^1x^m\dd x=\frac1{m+1},\qquad \int_0^1-x^m\log x\dd x=\frac1{(m+1)^2},
\]
the second by integration by parts (or by differentiating the first with respect to $m$).
Multiplying the two, with a second variable $y$, gives a cube:
\[
  \int_0^1\!\!\int_0^1-(xy)^m\log x\dd x\dd y=\frac1{(m+1)^3}.
\]
The roles of $x$ and $y$ are interchangeable, so we symmetrise: since $-\log x-\log y=-\log(xy)$,
\[
  \int_0^1\!\!\int_0^1-(xy)^m\log(xy)\dd x\dd y=\frac2{(m+1)^3}.
\]
Finally we sum over $m\ge0$, using $\sum_{m\ge0}(xy)^m=\frac1{1-xy}$ (all terms are nonnegative,
so Tonelli's theorem allows the interchange):
\begin{equation}\label{eq:zeta3int}
  \int_0^1\!\!\int_0^1\frac{-\log(xy)}{1-xy}\dd x\dd y=2\zeta(3).
\end{equation}

\subsection{Modifying the integral}

Having $\zeta(3)$ as an integral, we experiment: what happens if we put other functions in front of
the integrand? Inserting monomials $x^ry^s$ produces linear combinations of $\zeta(3)$ and $1$
(\cref{lem:moments}), and hence so does inserting any polynomial $P(x)P(y)$.

\begin{definition}
For $0<x,y<1$ let
\[
  K(x,y)=\frac{-\log(xy)}{1-xy}.
\]
For a polynomial $P$ put $J(P)=\int_0^1\!\int_0^1P(x)P(y)K(x,y)\dd x\dd y$, and for integers
$r,s\ge0$ put $M_{rs}=\int_0^1\!\int_0^1x^ry^sK(x,y)\dd x\dd y$.
\end{definition}

Write $H^{(q)}_j=\sum_{\ell=1}^j\ell^{-q}$, with $H^{(q)}_0=0$.

\begin{lemma}[Moments]\label{lem:moments}
$M_{rr}=2\bigl(\zeta(3)-H^{(3)}_r\bigr)$, and for $r\ne s$
\[
  M_{rs}=\frac{H^{(2)}_r-H^{(2)}_s}{r-s}\in\Q .
\]
In particular $\int\!\!\int K=M_{00}=2\zeta(3)$.
\end{lemma}

\begin{proof}
Expand $\frac1{1-xy}=\sum_{m\ge0}(xy)^m$ and use $-\log(xy)=-\log x-\log y$. All terms are
nonnegative, so Tonelli's theorem allows termwise integration:
\[
  M_{rs}=\sum_{m\ge0}\Bigl(\frac1{(m+r+1)^2(m+s+1)}+\frac1{(m+r+1)(m+s+1)^2}\Bigr).
\]
For $r=s$ this is $2\sum_{k>r}k^{-3}$: the sum in \eqref{eq:zeta3int} simply starts later. For $r>s$ put $d=r-s$ and $a=m+s+1$; the summand is
$\frac{2a+d}{a^2(a+d)^2}=\frac1d\bigl(\frac1{a^2}-\frac1{(a+d)^2}\bigr)$, and the sum telescopes to
$\frac1d\bigl(H^{(2)}_r-H^{(2)}_s\bigr)$. The case $s>r$ follows by symmetry.
\end{proof}

\begin{proposition}[Arithmetic]\label{prop:arith}
If $P(x)=\sum_{j=0}^nc_jx^j$ has integer coefficients, then
\[
  J(P)=2\Bigl(\sum_{j=0}^nc_j^2\Bigr)\zeta(3)-b(P),\qquad \Dn^3\,b(P)\in\Z .
\]
\end{proposition}

\begin{proof}
$J(P)=\sum_{r,s}c_rc_sM_{rs}$, and only the diagonal moments contain $\zeta(3)$. For
$r,s\le n$, each $j^3$ in $H^{(3)}_r$ divides $\Dn^3$, and each term of $M_{rs}$ with $r\ne s$ has
denominator $(r-s)j^2$ with $j\le n$, so $\Dn^3/((r-s)j^2)=\frac{\Dn}{r-s}\bigl(\frac{\Dn}{j}\bigr)^2\in\Z$.
\end{proof}

\begin{remark}
The denominator depends only on the \emph{degree} of $P$, not on the size of its coefficients.
For example $J(1-2x)=M_{00}-4M_{01}+4M_{11}=10\zeta(3)-12\approx0.0206$.
\end{remark}

% ======================================================================
\section{Designing the cancellation: Legendre polynomials}\label{sec:legendre}

We want integer polynomials $P,Q$ of degree $n$ for which
$\int_0^1\!\int_0^1P(x)Q(y)K(x,y)\dd x\dd y$ is small. We first solve the one-variable version:
make $\int_0^1P(x)f(x)\dd x$ small for a given function $f$. If $f$ has a Taylor expansion
$f(x)=a_0+a_1x+\dots+a_{n-1}x^{n-1}+(\text{higher order terms})$, we can make $P$ blind to the
low-order monomials by asking
\begin{equation}\label{eq:orth}
  \int_0^1x^jP(x)\dd x=0,\qquad j=0,\dots,n-1.
\end{equation}
Then $\int_0^1Pf$ sees only the higher-order part of $f$. These are $n$ linear equations in the
$n+1$ coefficients of $P$, so a nonzero solution of degree at most $n$ exists. It is unique up to
scale: the difference of two solutions with the same leading coefficient has degree below $n$ and
is orthogonal to itself, so it vanishes. Thus the requirement \emph{forces} $P$. An explicit
formula comes from integration by parts, since the $n$-th derivative kills every polynomial of
degree below $n$.

\begin{definition}
The \emph{shifted Legendre polynomial} of degree $n$ is
\[
  P_n(x)=\frac1{n!}\frac{\dd^n}{\dd x^n}\bigl[x^n(1-x)^n\bigr]
  =\sum_{k=0}^n(-1)^k\binom nk\binom{n+k}k x^k .
\]
\end{definition}

So $P_n$ has integer coefficients and $P_n(0)=1$; for example $P_1=1-2x$, $P_2=1-6x+6x^2$,
$P_3=1-12x+30x^2-20x^3$.

\begin{lemma}[Integration by parts]\label{lem:ibp}
For $f\in C^n[0,1]$,
\[
  \int_0^1P_n(x)f(x)\dd x=\frac{(-1)^n}{n!}\int_0^1x^n(1-x)^nf^{(n)}(x)\dd x .
\]
In particular $P_n$ satisfies \eqref{eq:orth}.
\end{lemma}

\begin{proof}
$x^n(1-x)^n$ has zeros of order $n$ at both endpoints, so all its derivatives of order below $n$
vanish there and every boundary term in $n$ integrations by parts is zero.
\end{proof}

The lemma is especially effective against reciprocal linear factors: for $t>0$,
$\frac{\dd^n}{\dd x^n}\frac1{x+t}=\frac{(-1)^nn!}{(x+t)^{n+1}}$, so
\begin{equation}\label{eq:recip}
  \int_0^1\frac{P_n(x)}{x+t}\dd x=\int_0^1\frac{x^n(1-x)^n}{(x+t)^{n+1}}\dd x>0 .
\end{equation}
An integral of an oscillating polynomial has become a positive integral with an explicit $n$-th
power.

We take $P=Q=P_n$.

\begin{definition}
$J_n=J(P_n)=\displaystyle\int_0^1\!\!\int_0^1P_n(x)P_n(y)\,\frac{-\log(xy)}{1-xy}\dd x\dd y.$
\end{definition}

By \cref{prop:arith}, $J_n$ satisfies \eqref{eq:arith}: $J_n=2a_n\zeta(3)-2b_n$ with
$a_n=\sum_k\binom nk^2\binom{n+k}k^2\in\Z$ and $\Dn^3b_n\in\Z$.

% ======================================================================
\section{Positivity}\label{sec:positivity}

The kernel $K(x,y)$ mixes the two variables (and has a logarithmic singularity at $0$), so the
one-variable identity \eqref{eq:recip} cannot be applied to it directly. We borrow a trick that
is common in differential equations: separation of variables, where a function of two variables is
written as a superposition of products, $u(x,t)=\sum v(x)w(t)$. Here we write $K$ as an integral,
over a parameter $t$, of products $v_t(x)w_t(y)$ in which each factor is a reciprocal linear
function, the best case for \eqref{eq:recip}.

\begin{lemma}[Separating the variables]\label{lem:sep}
For $0<x,y<1$,
\[
  K(x,y)=\int_0^\infty\frac{\dd t}{(x+t)(1+yt)} .
\]
\end{lemma}

\begin{proof}
For $0<a<b$, $\int_0^\infty\frac{\dd t}{(t+a)(t+b)}=\frac{\log b-\log a}{b-a}$ (partial fractions;
the logarithms cancel at infinity). Take $a=x$, $b=1/y$ and multiply numerator and denominator
by $y$.
\end{proof}

\begin{proposition}[Positive representation]\label{prop:pos}
\[
  J_n=\int_0^\infty\!\!\int_0^1\!\!\int_0^1
  \frac{\bigl[x(1-x)\,y(1-y)\,t\bigr]^n}{\bigl[(x+t)(1+yt)\bigr]^{n+1}}\dd x\dd y\dd t>0 .
\]
\end{proposition}

\begin{proof}
Since $|P_n|$ is bounded on $[0,1]$ and $\int\!\!\int\!\!\int\frac{\dd x\dd y\dd t}{(x+t)(1+yt)}=2\zeta(3)$,
Fubini's theorem lets us integrate in $x$ and $y$ first for each fixed $t>0$. By
\eqref{eq:recip}, the $x$-integral is $\int_0^1\frac{[x(1-x)]^n}{(x+t)^{n+1}}\dd x$, and since
$\frac{\dd^n}{\dd y^n}\frac1{1+yt}=\frac{(-1)^nn!\,t^n}{(1+yt)^{n+1}}$, the $y$-integral is
$t^n\int_0^1\frac{[y(1-y)]^n}{(1+yt)^{n+1}}\dd y$. The integrand is now nonnegative and not
identically zero.
\end{proof}

\begin{remark}
Bounding $|P_n(x)P_n(y)|$ at the start would throw away exactly the cancellation designed in
\cref{sec:legendre}. The positive representation must come first.
\end{remark}

% ======================================================================
\section{Exponential decay}\label{sec:decay}

Write $Q=(x+t)(1+yt)$ and $G=x(1-x)y(1-y)t/Q$, so $J_n=\int G^n/Q$.

\begin{lemma}\label{lem:max}
On $(0,1)^2\times(0,\infty)$ we have $G\le\rho:=(\sqrt2-1)^4=17-12\sqrt2$, with equality at
$x=y=\sqrt2-1$, $t=1$.
\end{lemma}

\begin{proof}
We have $J_n=\int G^n/Q$, so if $G\le\rho$ everywhere then $J_n\le\rho^n\int1/Q$. The task is to
maximise a function of three variables, and AM--GM reduces it to one. Put $u=\sqrt{xy}$ and write
\[
  G=\frac{xy\,(1-x)(1-y)}{Q/t},\qquad \frac Qt=\frac{(x+t)(1+yt)}t=1+xy+yt+\frac xt .
\]
\emph{Denominator.} By AM--GM, $yt+x/t\ge2\sqrt{xy}$, with equality when $t=\sqrt{x/y}$ balances
the two terms. Hence $Q/t\ge1+u^2+2u=(1+u)^2$.

\emph{Numerator.} For fixed $xy$, the product $(1-x)(1-y)=1-(x+y)+xy$ is largest when $x+y$ is
smallest, and $x+y\ge2\sqrt{xy}$. Hence $(1-x)(1-y)\le1-2u+u^2=(1-u)^2$.

Together, $G\le g(u)^2$ with $g(u)=\dfrac{u(1-u)}{1+u}$, $0<u<1$.

\emph{Finding the maximum.} The numerator of $g'(u)$ is $(1-2u)(1+u)-u(1-u)=1-2u-u^2$, which
vanishes at $u=\alpha:=\sqrt2-1$. Since $1-\alpha=\sqrt2\,\alpha$ and $1+\alpha=\sqrt2$, the
maximum value is $g(\alpha)=\alpha^2$.

\emph{Certifying it.} To prove $g(u)\le\alpha^2$ without calculus, consider the quadratic
$q(u)=\alpha^2(1+u)-u(1-u)$. If the claim is true, $q\ge0$ and $q(\alpha)=0$, so $\alpha$ is a
minimum of $q$ with value $0$: a double root. A monic quadratic with a double root at $\alpha$ is
$(u-\alpha)^2$, and indeed
\[
  q(u)=u^2-(1-\alpha^2)u+\alpha^2=(u-\alpha)^2,
\]
because $1-\alpha^2=2\alpha$, that is, $\alpha^2+2\alpha-1=0$, which is the equation that defined
$\alpha$. Therefore $g(u)\le\alpha^2$ and $G\le\alpha^4=(\sqrt2-1)^4=17-12\sqrt2$. All the
inequalities are equalities at $x=y=\alpha$, $t=\sqrt{x/y}=1$.
\end{proof}

\begin{theorem}\label{thm:decay}
For every $n\ge0$,
\[
  0<J_n\le2\zeta(3)\,(17-12\sqrt2)^n ,
\]
and in fact $\lim_{n\to\infty}J_n^{1/n}=17-12\sqrt2\approx0.0294$.
\end{theorem}

\begin{proof}
Positivity is \cref{prop:pos}. By \cref{lem:max}, $J_n\le\rho^n\int1/Q=\rho^n\int\!\!\int K=2\zeta(3)\rho^n$
(\cref{lem:sep,lem:moments}). For the limit, integrate only over a small box around
$(\alpha,\alpha,1)$, where $G>\rho-\varepsilon$ and $1/Q$ is bounded below.
\end{proof}

% ======================================================================
\section{Least common multiples without the prime number theorem}\label{sec:lcm}

We only need $\Dn^3$ to be smaller than $\rho^{-n}$ along \emph{some} sequence of indices.

\begin{lemma}\label{lem:factorial}
For every positive integer $N$, \ $N!=\prod_{k=1}^ND_{\lfloor N/k\rfloor}$.
\end{lemma}

\begin{proof}
Fix a prime $p$. The exponent of $p$ in $D_m$ is the number of $a\ge1$ with $p^a\le m$, so its
exponent on the right is $\sum_{a\ge1}\#\{k:\lfloor N/k\rfloor\ge p^a\}=\sum_{a\ge1}\lfloor N/p^a\rfloor$,
which is Legendre's formula for its exponent in $N!$.
\end{proof}

\begin{lemma}\label{lem:lcm}
For every real $C>e$ there are infinitely many $n$ with $\Dn<C^n$.
\end{lemma}

\begin{proof}
Suppose $D_m\ge C^m$ for all $m\ge M$, and let $L=\lfloor N/M\rfloor$. For $k\le L$ we have
$\lfloor N/k\rfloor\ge M$, and all other factors in \cref{lem:factorial} are at least $1$, so
\[
  N\log N\ge\log N!\ge\log C\sum_{k=1}^L\Bigl\lfloor\frac Nk\Bigr\rfloor\ge\log C\,(NH_L-L),
\]
where $H_L=\sum_{k\le L}1/k\ge\log(L+1)>\log N-\log M$. Dividing by $N$,
$\log N\ge\log C\,(\log N-\log M-1/M)$. Since $\log C>1$, this fails for large $N$.
\end{proof}

Since $e=\sum_k1/k!<1+\sum_{k\ge1}2^{1-k}=3$, we may take $C=3$: $\Dn<3^n$ for infinitely many $n$.

\begin{remark}
Beukers used the prime number theorem, $\Dn=e^{n+o(n)}$. Any $C$ with $e<C<\rho^{-1/3}\approx3.24$
would do here; the analytic estimate is what makes this interval nonempty.
\end{remark}

% ======================================================================
\section{The contradiction}\label{sec:contradiction}

\begin{theorem}[Ap\'ery]\label{thm:apery}
$\zeta(3)$ is irrational.
\end{theorem}

\begin{proof}
Suppose $\zeta(3)=p/q$ with $p,q\in\N$. By \cref{prop:arith}, $q\Dn^3J_n\in\Z$, and by
\cref{thm:decay} it is positive, hence at least $1$. For the infinitely many $n$ with $\Dn<3^n$,
\[
  1\le q\Dn^3J_n<2q\zeta(3)\bigl[27(17-12\sqrt2)\bigr]^n .
\]
Now $\rho=1/(17+12\sqrt2)<1/32$ because $\sqrt2>5/4$, so $27\rho<27/32<1$ and the right-hand side
tends to $0$. This is a contradiction.
\end{proof}

\begin{table}[ht]
\centering
\begin{tabular}{@{}ll@{}}
\toprule
Ingredient & What it provides\\
\midrule
Kernel $-\log(xy)/(1-xy)$ & only $\zeta(3)$ and rationals with denominator $\Dn^3$\\
Legendre polynomial $P_n$ & degree $n$, integer coefficients, $n$ vanishing moments\\
Separated kernel $+$ integration by parts & cancellation becomes a positive integral\\
Maximum of $G$ & the rate $(17-12\sqrt2)^n$\\
$N!=\prod D_{\lfloor N/k\rfloor}$ & infinitely many indices with $\Dn<3^n$\\
\bottomrule
\end{tabular}
\caption{Which ingredient supplies which property.}
\label{tab:ingredients}
\end{table}

% ======================================================================
\section{From Beukers' integral to Ap\'ery and Zudilin}\label{sec:zudilin}

The same numbers $J_n$ can be written as a sum of values of a rational function. This is the
language of Ap\'ery's original proof and of Zudilin's elementary proof~\cite{zudilin}, and the
starting point for deriving Ap\'ery's recurrence.

\subsection{The moment generating function of \texorpdfstring{$P_n$}{P\_n}}

Collect all the moments of $P_n$ in a single function (its Mellin transform):
\[
  T_n(s)=\int_0^1x^{s-1}P_n(x)\dd x,\qquad s>0 .
\]

\begin{proposition}\label{prop:mellin}
For $s>0$,
\[
  T_n(s)=(-1)^n\,\frac{(s-1)(s-2)\cdots(s-n)}{s(s+1)\cdots(s+n)} .
\]
\end{proposition}

\begin{proof}
Integrate by parts $n$ times as in \cref{lem:ibp}, with $f(x)=x^{s-1}$. Each boundary term at $0$ is
a product of a derivative of order $j<n$ of $x^n(1-x)^n$, which is $O(x^{n-j})$, and a derivative of
order $n-1-j$ of $x^{s-1}$, which is $O(x^{s-n+j})$; the product is $O(x^s)\to0$. The boundary terms
at $1$ vanish as before. Since $\frac{\dd^n}{\dd x^n}x^{s-1}=(s-1)(s-2)\cdots(s-n)\,x^{s-1-n}$,
\[
  T_n(s)=\frac{(-1)^n}{n!}(s-1)\cdots(s-n)\int_0^1x^{s-1}(1-x)^n\dd x
  =(-1)^n\frac{(s-1)\cdots(s-n)}{s(s+1)\cdots(s+n)},
\]
using the beta integral $\int_0^1x^{s-1}(1-x)^n\dd x=\frac{n!}{s(s+1)\cdots(s+n)}$ (by induction on
$n$, integrating by parts once).
\end{proof}

Every feature of this rational function has a meaning. The \emph{zeros} at $s=1,\dots,n$ are the
vanishing moments $T_n(j+1)=\int_0^1x^jP_n=0$ for $j<n$ (the condition \eqref{eq:orth}). The
\emph{poles} at $s=0,-1,\dots,-n$ come from integrating the monomials $x^k$, $k\le n$:
$T_n(s)=\sum_kc_k/(s+k)$.

\subsection{\texorpdfstring{$J_n$}{J\_n} as a sum of derivatives}

Since $P_n(x)P_n(y)$ appears in $J_n$, the transform appears twice. Put
\[
  R_n(s)=T_n(s)^2=\Bigl(\frac{(s-1)(s-2)\cdots(s-n)}{s(s+1)\cdots(s+n)}\Bigr)^2
  =\int_0^1\!\!\int_0^1(xy)^{s-1}P_n(x)P_n(y)\dd x\dd y .
\]

\begin{proposition}\label{prop:sum}
\[
  J_n=-\sum_{m=1}^\infty R_n'(m)=-\sum_{m=n+1}^\infty R_n'(m).
\]
\end{proposition}

\begin{proof}
Differentiating under the integral sign, with $\frac{\dd}{\dd s}(xy)^{s-1}=\log(xy)\,(xy)^{s-1}$,
\[
  -R_n'(s)=\int_0^1\!\!\int_0^1(xy)^{s-1}P_n(x)P_n(y)\bigl(-\log(xy)\bigr)\dd x\dd y .
\]
Now sum over $s=m+1$, $m\ge0$, exactly as in \cref{sec:kernel}: $\sum_{m\ge0}(xy)^m=\frac1{1-xy}$
turns the right-hand side into $J_n$. The interchange of sum and integral is justified because the
terms are dominated by $\|P_n\|_\infty^2\,(xy)^m(-\log(xy))$, whose sum $\|P_n\|_\infty^2K$ is
integrable by \eqref{eq:zeta3int}. Finally, $R_n$ has double zeros at $1,\dots,n$, so $R_n'$
vanishes there and the first $n$ terms are zero.
\end{proof}

The vanishing of the first $n$ terms explains part of the smallness of $J_n$, but not the
exponential rate: the remaining terms have both signs, and it is the positive integral of
\cref{prop:pos} that controls their cancellation. In the sum form, the same control comes instead
from Ap\'ery's recurrence, which is where Zudilin's proof goes next.

% ======================================================================
\begin{thebibliography}{9}
\bibitem{apery} R.~Ap\'ery, Irrationalit\'e de $\zeta(2)$ et $\zeta(3)$, \emph{Ast\'erisque}
  \textbf{61} (1979), 11--13.
\bibitem{beukers79} F.~Beukers, A note on the irrationality of $\zeta(2)$ and $\zeta(3)$,
  \emph{Bull.\ London Math.\ Soc.} \textbf{11} (1979), 268--272.
  \href{https://doi.org/10.1112/blms/11.3.268}{doi:10.1112/blms/11.3.268}.
\bibitem{beukers80} F.~Beukers, Legendre polynomials in irrationality proofs,
  \emph{Bull.\ Austral.\ Math.\ Soc.} \textbf{22} (1980), 431--438.
\bibitem{vdp} A.~van der Poorten, A proof that Euler missed\dots\ Ap\'ery's proof of the
  irrationality of $\zeta(3)$, \emph{Math.\ Intelligencer} \textbf{1} (1979), 195--203.
\bibitem{zudilin} W.~Zudilin, An elementary proof of Ap\'ery's theorem, 2002,
  \href{https://arxiv.org/abs/math/0202159}{arXiv:math/0202159}.
\end{thebibliography}

\end{document}
